\chapter{Stichwortverzeichnis (Index)}
\label{cha:Index}

Dieses Register enthält alle wichtigen grammatikalischen Begriffe und Themen.

\bigskip
\begin{center}
\begin{multicols}{2}
\textbf{A}
\item Akkusativ  \dotfill 21, 45
\item Aktiv  \dotfill 89
\item Artikel  \dotfill 31, 55
\item Attribut  \dotfill 78, 115

\textbf{B}
\item Bestimmter Artikel  \dotfill 31
\item B2-Niveau  \dotfill 95

\textbf{C}
\item C1-Niveau  \dotfill 103
\item C2-Niveau  \dotfill 111

\textbf{D}
\item Dativ  \dotfill 22, 47
\item Deklination  \dotfill 25
\item Demonstrativpronomen  \dotfill 37

\textbf{E}
\item E-Mail  \dotfill 121
\item Eigenname  \dotfill 33

\textbf{F}
\item Falsche Freunde  \dotfill 141
\item Femininum  \dotfill 26
\item FVG (Funktionsverbgefüge)  \dotfill 75

\textbf{G}
\item Genitiv  \dotfill 23, 52
\item Geschäftsbrief  \dotfill 119
\item Grammatik  \dotfill 1

\textbf{H}
\item Hauptsatz  \dotfill 65

\textbf{I}
\item Indirekte Rede  \dotfill 69
\item Indikativ  \dotfill 69
\item Infinitiv  \dotfill 12, 113

\textbf{K}
\item Kasus  \dotfill 21
\item Komparativ  \dotfill 91
\item Konjugation  \dotfill 9
\item Konjunktiv I  \dotfill 69
\item Konjunktiv II  \dotfill 71

\textbf{L}
\item Lokale Präpositionen  \dotfill 57

\textbf{M}
\item Maskulinum  \dotfill 26
\item Modalverb  \dotfill 15
\item Modalpartikel  \dotfill 99

\textbf{N}
\item Nebensatz  \dotfill 66
\item Negation  \dotfill 41
\item Neutrum  \dotfill 26
\item Nominalstil  \dotfill 87
\item Nomen  \dotfill 25

\textbf{O}
\item Objekt  \dotfill 21

\textbf{P}
\item Partizip I  \dotfill 14
\item Partizip II  \dotfill 13
\item Passiv  \dotfill 81
\item Perfekt  \dotfill 13
\item Personalpronomen  \dotfill 19
\item Plural  \dotfill 27
\item Possessivpronomen  \dotfill 38
\item Prädikat  \dotfill 9
\item Präsens  \dotfill 9
\item Präteritum  \dotfill 10
\item Präposition  \dotfill 55, 127
\item Pronomen  \dotfill 19

\textbf{R}
\item Reflexivpronomen  \dotfill 17
\item Relativpronomen  \dotfill 79
\item Relativsatz  \dotfill 79

\textbf{S}
\item Satzbau  \dotfill 63
\item SOV-Wortstellung  \dotfill 64
\item Subjekt  \dotfill 9
\item Subjunktor  \dotfill 66, 137
\item Superlativ  \dotfill 91

\textbf{T}
\item Temporale Präpositionen  \dotfill 58

\textbf{U}
\item Unbestimmter Artikel  \dotfill 31

\textbf{V}
\item Verbalstil  \dotfill 87
\item Verb  \dotfill 9
\item Verben mit Akkusativ  \dotfill 45, 155
\item Verben mit Dativ  \dotfill 47, 156
\item Verbposition  \dotfill 63

\textbf{W}
\item Wortstellung  \dotfill 63

\textbf{Z}
\item Zeitform  \dotfill 9
\item Zustandspassiv  \dotfill 84

\end{multicols}
\end{center}

\chapterend

\chapter{Glossar: Grammatikbegriffe}
\label{cha:Glossar}

\section{Deutsch-Englisch}

\bigskip
\begin{center}
\begin{tabular}{|l|l|}
\hline
Deutsch & English \\ \hline
Akkusativ & Accusative (case) \\ \hline
Adjektiv & Adjective \\ \hline
Adverb & Adverb \\ \hline
Aktiv & Active (voice) \\ \hline
Artikel & Article \\ \hline
Aussagesatz & Declarative sentence \\ \hline
Dativ & Dative (case) \\ \line
Deklination & Declension \\ \hline
Demonstrativpronomen & Demonstrative pronoun \\ \hline
Genitiv & Genitive (case) \\ \hline
Hauptsatz & Main clause \\ \hline
Indefinitpronomen & Indefinite pronoun \\ \hline
Indikativ & Indicative (mood) \\ \hline
Infinitiv & Infinitive \\ \hline
Konjunktion & Conjunction \\ \hline
Konjunktiv & Subjunctive (mood) \\ \hline
Maskulinum & Masculine (gender) \\ \hline
Modalverb & Modal verb \\ \hline
Nebensatz & Subordinate clause \\ \hline
Negation & Negation \\ \hline
Neutrum & Neuter (gender) \\ \hline
Nomen & Noun \\ \hline
Nominativ & Nominative (case) \\ \hline
Objekt & Object \\ \hline
Partizip & Participle \\ \hline
Passiv & Passive (voice) \\ \hline
Perfekt & Perfect (tense) \\ \hline
Personalpronomen & Personal pronoun \\ \hline
Plural & Plural \\ \hline
Possessivpronomen & Possessive pronoun \\ \hline
Prädikat & Predicate \\ \hline
Präsens & Present (tense) \\ \hline
Präteritum & Preterite (tense) \\ \hline
Präposition & Preposition \\ \hline
Pronomen & Pronoun \\ \hline
Reflexivpronomen & Reflexive pronoun \\ \hline
Relativpronomen & Relative pronoun \\ \hline
Relativsatz & Relative clause \\ \hline
Satz & Sentence \\ \hline
Singular & Singular \\ \hline
Subjekt & Subject \\ \hline
Subjunktor & Subordinating conjunction \\ \hline
Tempus & Tense \\ \hline
Verb & Verb \\ \hline
Wortstellung & Word order \\ \hline
\end{tabular}
\end{center}

\section{Begriffe erklärt}

\bigskip
\noindent\textbf{Kasus (Fall):} Die grammatikalische Funktion eines Nomens oder Pronomens im Satz.

\bigskip
\noindent\textbf{Nominativ:} Subjektform. „Der Mann kommt." (Who?)

\bigskip
\noindent\textbf{Akkusativ:} Objektform (direkt). „Ich sehe den Mann." (Whom?)

\bigskip
\noindent\textbf{Dativ:} Objektform (indirekt). „Ich helfe dem Mann." (To whom?)

\bigskip
\noindent\textbf{Genitiv:} Besitzform. „Das Auto des Mannes." (Whose?)

\bigskip
\noindent\textbf{Modus (Stimmung):} Indikativ = real, Konjunktiv = hypothetisch/irreal.

\bigskip
\noindent\textbf{Genus (Geschlecht):} Maskulinum (der), Femininum (die), Neutrum (das).

\bigskip
\noindent\textbf{Diathese (Stimme):} Aktiv vs. Passiv.

\bigskip
\noindent\textbf{Tempus (Zeit):} Präsens, Präteritum, Perfekt, Plusquamperfekt, Futur I/II.

\chapterend